%% ma: L11-B16-0067
%% lop: 11
%% chuong: 5
%% bai: 16
%% dang: 11.16.2
%% loai: TLN
%% muc: VD
%% dap_an: "9"
%% nguon: "tai-lieu-hoc-tap-mon-toan-lop-11-theo-chuong-trinh-giao-duc-pho-thong-2018.pdf (Lớp Toán Thầy Duy, Phạm Lê Duy, THPT Chu Văn An), Chương 3 Bài 2, Phần 3, Câu 4 (tr.381)"
%% kien_thuc: "Lớp 11 Bài 16 (giới hạn của hàm số)"
%% trang_thai: cho-duyet
%% ly_do: "Dạng mới đề xuất (11.16.x), chờ giáo viên duyệt dạng; tài liệu gốc không có lời giải mẫu, đáp án do kho tự giải và kiểm bằng sympy."
\begin{ex}
Biết rằng giới hạn \[\lim\limits_{x\to0}\dfrac{\sqrt{x+9}+\sqrt{x+16}-7}x=\lim\limits_{x\to0}\left[\dfrac a{\sqrt{x+9}+b}+\dfrac c{\sqrt{x+16}+d}\right]\] với $a,b,c,d$ là các số nguyên dương. Tính tổng các số $a,b,c,d$.
\shortans{9}
\loigiai{\sloppy \[\dfrac{\sqrt{x+9}+\sqrt{x+16}-7}x=\dfrac{\sqrt{x+9}-3}x+\dfrac{\sqrt{x+16}-4}x=\dfrac1{\sqrt{x+9}+3}+\dfrac1{\sqrt{x+16}+4}.\] Vậy $a=1$, $b=3$, $c=1$, $d=4$ và tổng $=9$.}
\end{ex}
